%% AMS-LaTeX Created with the Wolfram Language : www.wolfram.com

\documentclass[12pt]{article}
\usepackage[letterpaper,margin=1in]{geometry}
\usepackage{amsmath, amssymb, graphics, setspace}
\usepackage{graphicx} 
\graphicspath{ {./figs/} }
\usepackage{tcolorbox}
\usepackage[framed,numbered,autolinebreaks,useliterate]{mcode} % Removed this line
\newcommand{\mathsym}[1]{{}}
\newcommand{\unicode}[1]{{}}

\usepackage[english]{babel}
\newtheorem{theorem}{Theorem}

\usepackage{caption}
\usepackage{subcaption}

\usepackage{tabulary}
\usepackage{multirow}
\usepackage{lscape}

\usepackage{hyperref}
\hypersetup{
    colorlinks=true,
    linkcolor=blue,
    filecolor=magenta,      
    urlcolor=cyan,
    pdftitle={Overleaf Example},
    pdfpagemode=FullScreen,
    }

\def\mathbi#1{\textbf{\em #1}}

\begin{document}
\doublespacing	

\title{CE 401/5501: Applied AI \\ \large Topic 16: Symbolic AI toward Next-gen AI for Engineering}
\author{Dr. ZhiQiang Chen}
\date{}
\maketitle


% Section 1: Revisiting Good-Old-Fashioned AI (GOFAI)
\section{Revisiting Symbolic AI: Foundations and Classic Methods}
\label{sec:gofai}
Symbolic AI, or GOFAI, dominated early AI research by emphasizing explicit knowledge representation and logic-based reasoning. Key methods include:

\subsection{Logic Systems}
First-order logic (FOL) formalizes domain knowledge with predicates, quantifiers, and rules. For example, in engineering design, material compatibility can be encoded as:
\begin{align}
\forall x \, (\text{Material}(x) \rightarrow \exists y \, (\text{StressThreshold}(x, y)))
\end{align}
Seminal work by McCarthy (1958) introduced FOL for AI, while Robinson’s \textit{resolution refutation} (1965) enabled automated theorem proving:
\begin{align}
\frac{\neg P \vee Q, \quad \neg Q \vee R}{\neg P \vee R} \quad \text{(Resolution Rule)}
\end{align}

\subsection{Expert Systems}
MYCIN (Shortliffe, 1976) demonstrated rule-based reasoning for medical diagnosis, inspiring engineering systems like XCON for configuring computer systems. Rules follow:
\begin{align}
\text{IF } \text{Condition}_1 \land \text{Condition}_2 \rightarrow \text{Action}
\end{align}

\subsection{Search Algorithms}
A* search (Hart et al., 1968) optimizes paths using heuristics \( h(n) \):
\begin{align}
f(n) = g(n) + h(n)
\end{align}
where \( g(n) \) is the path cost and \( h(n) \) is admissible (never overestimates).

% Section 2: Relevance to Modern GenAI
\section{Symbolic AI in the Generative AI Era}
\label{sec:genai}
Modern GenAI (e.g., LLMs) struggles with structured reasoning, necessitating hybrid approaches:

\subsection{Bridging Hallucinations with Knowledge Graphs (KGs)}
KGs ground LLMs in factual triples \((h, r, t)\) (e.g., \textit{(Steel, hasDensity, 8,050 kg/m³)}). Google’s Gemini uses KGs to reduce hallucinations by 40\% in technical domains.

\subsection{Constraint Solvers for Logical Consistency}
Neuro-symbolic systems like Microsoft’s Minerva integrate transformers with satisfiability modulo theories (SMT):
\begin{align}
\text{SMT-Lib: } (declare-fun x () Int) \quad (assert (> x 0))
\end{align}

% Section 3: AI for 21st Century Engineering
\section{Engineering Demands Hybrid AI}
\label{sec:engineering}
Modern engineering requires integrating physical laws, simulations, and collaboration:

\subsection{Physics-Informed Neural Networks (PINNs)}
Raissi et al. (2019) enforce PDE constraints on ML models:
\begin{align}
\mathcal{L} = \underbrace{\frac{1}{N} \sum_{i=1}^N |y_i - \hat{y}_i|^2}_{\text{Data Loss}} + \lambda \underbrace{|\nabla \hat{y} - f(\hat{y})|^2}_{\text{Physics Loss}}
\end{align}

\subsection{Multi-Agent Systems for Collaboration}
Distributed constraint optimization (DCOP) coordinates teams:
\begin{align}
\max \sum_{a_i \in A} U_i(a_i), \quad \text{s.t. } C_{ij}(a_i, a_j) \leq 0
\end{align}

% Section 4: Neuro-Symbolic AI
\section{Neuro-Symbolic AI: Architecture and Impact}
\label{sec:neurosym}
Neuro-symbolic AI unifies neural networks (pattern recognition) and symbolic systems (reasoning):

\subsection{Architectural Frameworks}
\begin{itemize}
\item \textbf{Neural-Guided Symbolic Reasoning}: LLMs propose candidate solutions, validated by symbolic solvers (e.g., Codex + Python interpreter).
\item \textbf{Symbolic Knowledge Distillation}: Extract rules from neural models (e.g., DeepMind’s Tracr compiles programs to neural weights).
\end{itemize}

\subsection{Logical Neural Networks (LNNs)}
IBM’s LNNs (Riegel et al., 2020) unify logic and gradients:
\begin{align}
\sigma(w \cdot \text{AND}(x, y) + b) \quad \text{(Differentiable Logic Gates)}
\end{align}

% Section 5: Top 3 Symbolic AI Methods
\section{Top 3 Symbolic AI Methods for Engineering}
\label{sec:top3}
\begin{enumerate}
\item \textbf{Knowledge Graphs (KGs) with Graph Neural Networks} \\
Integrate KGs (e.g., Neo4j) with GNNs for relational reasoning:
\begin{align}
h_v^{(k)} = \sigma\left(W^{(k)} \cdot \text{AGGREGATE}(\{h_u^{(k-1)}, \forall u \in \mathcal{N}(v)\})\right)
\end{align}

\item \textbf{Constraint Programming (CP)} \\
Solve NP-hard scheduling via MiniZinc:
\begin{align}
% \text{var 1..n: x; constraint alldifferent([x_1, ..., x_n]);}
\end{align}

\item \textbf{Automated Planning with STRIPS} \\
Define engineering workflows as \( \langle P, A, I, G \rangle \):
\begin{align}
\text{Action: } \text{Preconditions} \rightarrow \text{Effects}
\end{align}
\end{enumerate}

% Summary
\section{Summary}
Hybrid neuro-symbolic AI is critical for engineering’s future, combining GenAI’s adaptability with symbolic rigor. Prioritize KGs, constraint solvers, and automated planning to address safety, scalability, and collaboration.

% References
\bibliographystyle{unsrt}
\bibliography{topic16-references}

\end{document}